\section{Introducción}

Este informe documenta el Trabajo práctico 1 del curso de Principios de
Inteligencia Artificial. El trabajo consiste en
implementar en PyTorch \cite{paszke2019pytorch} tres algoritmos de
optimización (descenso del gradiente, RMSProp y enjambre de partículas) y
compararlos sobre tres funciones de prueba:

\begin{align*}
f_0(x,y) &= x^2+y^2, \\
f_1(x,y) &= \text{función de Ackley \cite{ackley1987connectionist}}, \\
f_2(x,y) &= \text{función de Himmelblau \cite{himmelblau1972applied}},
\end{align*}

definidas sobre el dominio $x,y\in[-10,10]$. Las demostraciones iniciales y
la calibración de hiperparámetros usan $P=25$ iteraciones; la evaluación
final usa un máximo de 50 iteraciones, como pide el enunciado. Todos los
puntos iniciales se muestrean uniformemente en $[-10,10]^2$.

La implementación se presenta en el notebook principal, acompañado por el
módulo auxiliar \texttt{utils.py}, que reúne las funciones de análisis,
visualización y presentación de tablas. Todos los números, tablas y figuras
de este informe provienen de una ejecución completa del notebook desde un
kernel limpio; las figuras se exportan automáticamente durante esa ejecución.

\subsection{Protocolo común}

\begin{itemize}
    \item \textbf{Criterio de convergencia.} Una corrida converge en la
    primera iteración $t$ que cumple $|f(x_t)-f^*|\leq10^{-3}$, con
    $f^*=0$ en las tres funciones. En PSO se evalúa el mejor valor global
    del enjambre en cada iteración. La iteración $t=0$ corresponde al punto
    inicial.
    \item \textbf{Calibración.} Se usa Optuna \cite{akiba2019optuna} con el
    muestreador TPE, 25 ensayos por estudio, 25 iteraciones fijas por
    corrida y los mismos 6 puntos iniciales para todos los estudios
    (generados con \texttt{random.Random(42)}). El objetivo que se minimiza
    es el promedio de $f$ al final de las 6 corridas, sin penalizaciones.
    Un ensayo que diverge se poda (\texttt{TrialPruned}).
    \item \textbf{Evaluación.} Con los mejores hiperparámetros se ejecutan
    10 corridas de 50 iteraciones desde los mismos 10 puntos iniciales
    (generados con \texttt{random.Random(20260920)}) para los tres
    algoritmos. En PSO, cada punto es la posición inicial de la primera
    partícula.
    \item \textbf{Reproducibilidad.} Las demostraciones usan puntos generados
    con \texttt{random.Random(2026)} y las partes estocásticas de PSO usan
    \texttt{torch.manual\_seed}, por lo que una nueva ejecución del notebook en
    el mismo entorno reproduce exactamente los resultados aquí reportados
    (ver las limitaciones en la Sección~\ref{sec:limitaciones}).
\end{itemize}
